н\documentclass[a4paper,14pt]{extarticle}  % extarticle поддерживает 14pt

% Кодировка и язык
\usepackage[utf8]{inputenc}   
\usepackage[T1]{fontenc}      
\usepackage[russian]{babel}   
\usepackage[colorlinks=true, urlcolor=blue]{hyperref}
% Шрифты
\usepackage{titlesec}
\usepackage{xcolor}

\usepackage{titlesec}
\usepackage{xcolor}
\titleformat{\section}
  {\bfseries\color{black}\LARGE}{\thesection}{1em}{}

\titleformat{\subsection}
  {\bfseries\color{black}\Large}{\thesubsection}{1em}{}
\usepackage{mathptmx}         

% Поля, таблицы, картинки
\usepackage{float}
\usepackage{geometry}         
\usepackage{graphicx}         
\usepackage{array}            
\usepackage{booktabs}         
\usepackage{caption}          
\usepackage{amsmath,amssymb}  
\setcounter{secnumdepth}{0}
% Настройки полей страницы 
\usepackage{geometry}
\geometry{
  top=2cm,
  bottom=2cm,
  left=3cm,
  right=2cm
}

% Абзацные отступы и межстрочный интервал
\setlength{\parindent}{1.25cm}
\renewcommand{\baselinestretch}{1.5}

% Заголовок
\title{Solutions of Diophanite Equation $x^3+y^3=x^2+kxy+y^2$}
\author{Eric Maxud}
\date{} % можно поставить дату или оставить пустой

\begin{document}
\maketitle

\section{Abstract}
First time We saw problem of similar form while listing AOPS. The original statement of the problem was as follows:
\\
\\
\noindent\fbox{\begin{minipage}{\textwidth}
\paragraph{BMO 2017/1} Find all ordered pairs of positive integers$ (x, y)$ such that:$$x^3+y^3=x^2+42xy+y^2.$$
\end{minipage}}
\\
\\
After that, We became interested in studying solutions to the generalized case. Thus we study integer solutions of equation $x^3+y^3=x^2+kxy+y^2$
\section{Introduction}
For small cases:
\begin{enumerate}
    \item if $k=0$, $$x^3+y^3=x^2+y^2$$ $$x^2-xy+y^2|x^2+y^2$$ Then $$x^2-xy+y^2|xy$$ So $x=y=0$ and $x=y=1$ are only solutions
    \item if $k=-1$, similarly we can get answers are all $x,y$ such that $x+y=1$ 
    \item if $k=1$, $$x^3+y^3=x^2+xy+y^2$$ Let
\[
s = x + y, \quad p = xy.
\]
Then
\[
x^3 + y^3 = (x+y)^3 - 3xy(x+y) = s^3 - 3ps,
\]
\[
x^2 + xy + y^2 = (x+y)^2 - xy = s^2 - p.
\]

Substituting, we get
\[
s^3 - 3ps = s^2 - p.
\]
Rearranging:
\[
s^3 - s^2 - 3ps + p = 0,
\]
\[
s^2(s - 1) - p(3s - 1) = 0.
\]
Thus,
\[
p(3s - 1) = s^2(s - 1).
\]

Observe that
\[
\gcd(3s - 1, s) = 1,
\]
since any common divisor divides
\[
(3s - 1) - 3s = -1.
\]

Also,
\[
\gcd(3s - 1, s - 1) \mid ((3s - 1) - 3(s - 1)) = 2.
\]

Hence from
\[
3s - 1 \mid s^2(s - 1)
\]
it follows that
\[
3s - 1 \mid 2.
\]

Therefore,
\[
3s - 1 \in \{-2, -1, 1, 2\}.
\]

Checking integer values:
\[
3s - 1 = -1 \Rightarrow s = 0, \quad
3s - 1 = 2 \Rightarrow s = 1.
\]

\textbf{Case 1:} $s = 0$  
\[
p(-1) = 0 \Rightarrow p = 0.
\]
Thus,
\[
x + y = 0, \quad xy = 0 \Rightarrow x = y = 0.
\]

\textbf{Case 2:} $s = 1$  
\[
p(2) = 0 \Rightarrow p = 0.
\]
Thus,
\[
x + y = 1, \quad xy = 0,
\]
so
\[
(x,y) = (1,0), \ (0,1).
\]

\textbf{Answer:}
\[
(x,y) = (0,0), \ (1,0), \ (0,1).
\]
\item
if \(k=2\), then the equation becomes
\[
x^3+y^3=x^2+2xy+y^2=(x+y)^2,
\]
and every pair with
\[
x+y=0
\]
is a solution, so there are infinitely many integer solutions.
\end{enumerate}
\space
\section{Theorems and their proves.}
Despite the small number of solutions for small cases, an equation of this type has quite a lot of solutions.
\\
\\
\noindent\fbox{\begin{minipage}{\textwidth}
\paragraph{Theorem 1:} For every positive integers $n,l$, there is exist $k>l$, such that $$x^3+y^3=x^2+kxy+y^2$$ has more than $n$ solutions in $\mathbb{Z}$.
\end{minipage}}
\\
\\
\textbf{Proof.}
We look for solutions of the form
\[
y=mx
\]
for some positive integer \(m\). Substituting into
\[
x^3+y^3=x^2+kxy+y^2
\]
gives
\[
x^3+(mx)^3=x^2+kx(mx)+(mx)^2,
\]
hence
\[
x^3(1+m^3)=x^2(1+km+m^2).
\]
Therefore, if
\[
1+m^3 \mid 1+km+m^2,
\]
then
\[
x=\frac{1+km+m^2}{1+m^3},\qquad y=mx
\]
is an integer solution.

So for each chosen \(m\), it is enough to require
\[
1+km+m^2\equiv 0 \pmod{1+m^3}.
\]
A sufficient condition is
\[
k\equiv m(m-1)\pmod{1+m^3},
\]
because then
\[
1+km+m^2 \equiv 1+m^2(m-1)+m^2 = 1+m^3 \equiv 0 \pmod{1+m^3}.
\]

Now choose
\[
m_i=2^{2^i}\qquad (i=1,2,\dots,n+1),
\]
and let
\[
q_i=1+m_i^3=1+2^{3\cdot 2^i}.
\]

We claim that the integers \(q_1,q_2,\dots,q_{n+1}\) are pairwise coprime.  
Indeed, suppose \(i<j\) and some integer \(d>1\) divides both \(q_i\) and \(q_j\). Then
\[
2^{3\cdot 2^i}\equiv -1 \pmod d.
\]
Raising both sides to the power \(2^{j-i}\), we obtain
\[
2^{3\cdot 2^j}\equiv (-1)^{2^{j-i}}=1 \pmod d.
\]
On the other hand, since \(d\mid q_j\), we also have
\[
2^{3\cdot 2^j}\equiv -1 \pmod d.
\]
This is impossible. Hence
\[
\gcd(q_i,q_j)=1 \qquad \text{for all } i\ne j.
\]

Therefore, by the Chinese Remainder Theorem, there exists an integer \(k\) such that
\[
k\equiv m_i(m_i-1)\pmod{q_i}
\qquad (i=1,2,\dots,n+1).
\]
Moreover, all such integers form an arithmetic progression, so we may choose one with
\[
k>l.
\]

For each \(i\), define
\[
x_i=\frac{1+km_i+m_i^2}{1+m_i^3},\qquad y_i=m_i x_i.
\]
By construction, \(x_i\in\mathbb Z\), hence \((x_i,y_i)\in\mathbb Z^2\). Also,
\[
x_i^3+y_i^3=x_i^3(1+m_i^3)
= x_i^2(1+km_i+m_i^2)
= x_i^2+kx_i y_i+y_i^2,
\]
so \((x_i,y_i)\) is a solution of the equation.

These solutions are distinct, since
\[
\frac{y_i}{x_i}=m_i
\]
and the numbers \(m_1,m_2,\dots,m_{n+1}\) are pairwise distinct.

Thus the equation has at least \(n+1\) distinct integer solutions. In particular, it has more than \(n\) solutions in \(\mathbb Z\).

This proves the Theorem 1.
\\
\\
\noindent\fbox{\begin{minipage}{\textwidth}
\paragraph{Theorem 2:} Let
\[
N_k=\#\{(x,y)\in \mathbb Z^2:\ x^3+y^3=x^2+kxy+y^2\}.
\]
If \(k\neq -1,2\), then
\[
N_k\le 4\,\tau\!\left(\left|(k-2)^2(k+1)\right|\right),
\]
where \(\tau(m)\) denotes the number of positive divisors of \(m\).

Moreover, for \(k=-1\) and \(k=2\), the equation has infinitely many integer solutions.
\end{minipage}}
\\
\\
\textbf{Proof.}
Consider the equation
\[
x^3+y^3=x^2+kxy+y^2.
\]
Set
\[
s=x+y,\qquad p=xy.
\]
Then
\[
x^3+y^3=(x+y)^3-3xy(x+y)=s^3-3ps,
\]
and
\[
x^2+kxy+y^2=(x+y)^2+(k-2)xy=s^2+(k-2)p.
\]
Hence
\[
s^3-3ps=s^2+(k-2)p,
\]
so
\[
s^3-s^2=p(3s+k-2),
\]
that is,
\[
s^2(s-1)=p(3s+k-2).
\]
Let
\[
D=3s+k-2.
\]
Then
\[
D\mid s^2(s-1).
\]

We first note that if \(D=0\), then
\[
s^2(s-1)=0,
\]
so \(s=0\) or \(s=1\). Since \(D=3s+k-2\), this gives \(k=2\) or \(k=-1\). Therefore, for
\[
k\neq -1,2,
\]
we must have
\[
D\neq 0.
\]

We now prove that
\[
|D|\mid |(k-2)^2(k+1)|.
\]

Let \(q^a\parallel D\), where \(q\) is a prime. Since
\[
D\mid s^2(s-1)
\]
and
\[
\gcd(s,s-1)=1,
\]
the prime \(q\) divides exactly one of \(s\) and \(s-1\).

\medskip

\textit{Case 1: \(q\mid s-1\).}

Since \(q\nmid s\), from \(q^a\mid s^2(s-1)\) we get
\[
q^a\mid (s-1).
\]
But also \(q^a\mid D\), hence
\[
q^a\mid D-3(s-1)=k+1.
\]

\medskip

\textit{Case 2: \(q\mid s\).}

Write
\[
q^b\parallel s.
\]
Since \(q\nmid s-1\), from \(q^a\mid s^2(s-1)\) we obtain
\[
a\le 2b.
\]
Also, because \(q^a\mid D\) and \(q^b\mid s\), we have
\[
q^{\min(a,b)}\mid \gcd(D,s).
\]
But
\[
D-3s=k-2,
\]
so
\[
q^{\min(a,b)}\mid (k-2).
\]
If \(a\le b\), then \(q^a\mid (k-2)\). If \(a>b\), then \(q^b\mid (k-2)\), and since \(a\le 2b\), it follows that
\[
q^a\mid (k-2)^2.
\]
Thus in all cases,
\[
q^a\mid (k-2)^2.
\]

Therefore, every prime power divisor of \(D\) divides either \(k+1\) or \((k-2)^2\), and so
\[
|D|\mid |(k-2)^2(k+1)|.
\]

Now let
\[
M=(k-2)^2(k+1).
\]
Every integer solution \((x,y)\) gives rise to a divisor \(D\) of \(M\), and \(s=x+y\) is then uniquely determined by
\[
s=\frac{D-k+2}{3}.
\]
After that, \(p=xy\) is uniquely determined from
\[
p=\frac{s^2(s-1)}{D}.
\]

Finally, \(x\) and \(y\) are the roots of
\[
t^2-st+p=0.
\]
For fixed \(s\) and \(p\), this quadratic gives at most two ordered integer pairs \((x,y)\), namely \((x,y)\) and \((y,x)\).

Since \(M\neq 0\) for \(k\neq -1,2\), the number of integer divisors of \(M\) is
\[
2\tau(|M|).
\]
Hence
\[
N_k\le 2\cdot 2\tau(|M|)=4\tau(|M|).
\]
Therefore
\[
N_k\le 4\,\tau\!\left(\left|(k-2)^2(k+1)\right|\right).
\]

\bigskip

It remains to note the exceptional cases from introduction.
This completes the proof.
\section{Numericial examples}
\[
\text{Below, all solution sets are sets of ordered pairs }(x,y)\in\mathbb Z^2.
\]

\begin{center}
\renewcommand{\arraystretch}{1.25}
\begin{tabular}{c l}
\toprule
\(\boldsymbol{k}\) & \(\boldsymbol{\text{Integer solutions of }x^3+y^3=x^2+kxy+y^2}\) \\
\midrule
\(-5\) & \((0,0),\ (1,0),\ (0,1)\) \\

\(-4\) & \((0,0),\ (1,0),\ (0,1),\ (-1,-1)\) \\

\(-1\) & infinitely many: \((x,y)=(t,1-t)\), \(t\in\mathbb Z\) \\

\(0\) & \((0,0),\ (1,0),\ (0,1),\ (1,1)\) \\

\(1\) & \((0,0),\ (1,0),\ (0,1)\) \\

\(2\) & infinitely many: \((x,y)=(t,-t)\), \(t\in\mathbb Z\) \\

\(4\) & \((0,0),\ (1,0),\ (0,1),\ (3,3)\) \\

\(6\) & \((0,0),\ (1,0),\ (0,1),\ (-2,1),\ (1,-2),\ (1,3),\ (3,1),\ (4,4)\) \\
\bottomrule
\end{tabular}
\end{center}
\section{Conclusion}
For an equation of this type, it is most likely impossible to derive a general form of the solution, especially from previously proven facts, for almost any $k$, there is a finite number of solutions, but for specific $k$ there can be a lot of them.
\section{References}
\begin{enumerate}
    \item \href{https://link.springer.com/book/10.1007/978-1-4757-1920-8}{The Arithmetic of Elliptic Curves} Joseph H. Silverman

    \item \href{https://link.springer.com/book/10.1007/978-1-4757-4252-7}{Rational Points on Elliptic Curves} Joseph H. Silverman, John Tate

    \item \href{https://www.cambridge.org/core/books/algorithmic-resolution-of-diophantine-equations/8140C8D3E91A8C933D21976163F3D905}{The Algorithmic Resolution of Diophantine Equations: A Computational Cookbook} Nigel P. Smart

    \item \href{https://link.springer.com/series/7604}{Number Theory. Volume I: Tools and Diophantine Equations} Henri Cohen

    \item \href{https://www.sciencedirect.com/science/article/pii/0022314X89900140}{On the Practical Solution of the Thue Equation} N. Tzanakis, B. M. M. de Weger

    \item \href{https://webusers.imj-prg.fr/~michel.waldschmidt/articles/pdf/ProcHRI2017ThueEquations.pdf}{Thue Diophantine Equations} Michel Waldschmidt

    \item \href{https://eudml.org/doc/144332}{The Number of Solutions of Decomposable Form Equations} Jan-Hendrik Evertse

    \item \href{https://dms.umontreal.ca/~andrew/Courses/RationalPtsOnCurves.pdf}{Rational and Integral Points on Curves} Andrew Granville
\end{enumerate}
\end{document}
